\section{Conclusion}

We built the dataset that makes crack re-identification testable ---
\textsc{CrackID}, 140 photographs over 17 walls with 212 identities, its
annotation provenance ($45\%$ merged) stated rather than assumed --- and the
benchmark's main finding is negative and structural: on the decision a record
requires, nine matchers across four families and two orders of compute buy
between $0.000$ and $0.054$ over chance, and only LoFTR survives a
validation-calibrated threshold. Ordering a gallery and deciding an identity are
not the same capability. The difficulty is not crack appearance: a geometric
reference baseline that registers on wall texture with cracks masked out and
Chamfer-matches projected masks reaches $\mathrm{DIR}@\mathrm{FAR}{=}0.1$ of
$0.768$ against $0.404$ for the best context method while answering $48\%$ of
pairs --- and a control exploiting that alignment coverage without examining a
crack scores $\mathbf{0.000}$. The open-set advantage has a specific, unshared
source: an aligned pair supports a \emph{licensed negative}, positive evidence
of absence, where a crop embedding can only report low similarity. Two bounds
remain addressable --- a real second capture session turns the benchmark from
viewpoint to cross-visit, and the annotation audit awaits its verdicts. The
$48\%$ coverage is also a deployment gap in its own right, not merely a scoring
issue: smooth-plaster walls can be categorically unanswerable. Finally, the
synthetic revisit protocol reverses the headline ranking and open-set result in
favour of OSNet$+$context, while the geometric pipeline retains only a narrow
assignment-$F_1$ advantage; the present paper therefore establishes a promising
within-visit geometric reference, not a solved multi-visit maintenance system.